\par\indent\SourceLeaf{236}%
III. Setzt man \(x=r\cos\varphi,\ y=r\sin\varphi\), so erhält man
\[
z=\tfrac{1}{2}(T\cos\varphi^2+V\sin\varphi^2)rr+\Omega
\]
Daraus folgt, dass ein Schnitt mit einer Ebene, die in \(A\) normal zur Fläche ist und mit der Achse \(x\) den Winkel \(\varphi\) bildet, eine ebene Kurve mit folgendem Krümmungsradius im Punkt \(A\) ergibt
\[
=\frac{1}{T\cos\varphi^2+V\sin\varphi^2}
\]
\par\indent IV. Gilt also \(T=V\), so sind die Krümmungsradien in \textit{allen} Normalebenen gleich. Sind \(T\) und \(V\) dagegen ungleich, so liegt \(T\cos\varphi^2+V\sin\varphi^2\) für jeden Wert des Winkels \(\varphi\) zwischen \(T\) und \(V\). Daher gehören die Krümmungsradien der Hauptschnitte, die in I und II betrachtet wurden, offenbar zu den extremen Krümmungen. Haben \(T\) und \(V\) dasselbe Vorzeichen, so gehört der eine zur größten, der andere zur kleinsten Krümmung. Haben \(T\) und \(V\) hingegen entgegengesetzte Vorzeichen, so gehört der eine zur größten Konvexität, der andere zur größten Konkavität. Diese Folgerungen enthalten fast alles, was der berühmte \textsc{Euler} zuerst über die Krümmung gekrümmter Flächen gelehrt hat.

V. Das Krümmungsmaß der gekrümmten Fläche im Punkt \(A\) erhält die sehr einfache Form \(k=TV\). Daraus folgt der folgende Satz.

\textsc{Satz.} \textit{Das Krümmungsmaß in einem beliebigen Punkt der Fläche ist gleich einem Bruch mit dem Zähler eins und dem Produkt der beiden extremen Krümmungsradien in Schnitten durch Normalebenen als Nenner.}

Zugleich ist klar, dass das Krümmungsmaß bei konkav-konkaven oder konvex-konvexen Flächen positiv ist (der Unterschied ist unwesentlich), bei konkav-konvexen Flächen dagegen negativ. Besteht eine Fläche aus Teilen beider Arten, so muss das Krümmungsmaß an den Grenzen zwischen ihnen verschwinden. Auf die Beschaffenheit gekrümmter Flächen, deren Krümmungsmaß überall verschwindet, gehen wir weiter unten ausführlicher ein.

\Article{9}
Die am Ende von Artikel 7 angegebene allgemeine Formel für das Krümmungsmaß ist die einfachste von allen, da sie nur fünf Größen enthält. Zu einer komplizierteren Formel mit neun Größen gelangen wir, wenn wir
